\documentclass[12pt]{article}
% Préambule commun à tous les documents extraits de « La Revue CAPES »
\usepackage[T1]{fontenc}
\usepackage[utf8]{inputenc}
\usepackage{lmodern}
\usepackage{amsmath,amssymb,amsthm,amsfonts,mathrsfs}
\usepackage{setspace}
\usepackage{listings}
\usepackage{pgfplots}
\pgfplotsset{compat=1.18}
\usepackage{longtable}
\usepackage{adjustbox}
\usepackage{lettrine}
\usepackage{tkz-tab}
\usepackage{changepage}
\usepackage{pdflscape}
\usepackage{graphicx}
\usepackage{tikz}
\usepackage{xcolor}
\usepackage[a4paper,top=2cm,bottom=2.5cm,left=2.5cm,right=2cm]{geometry}
\usepackage{fancyhdr}
\usepackage{draftwatermark}
\SetWatermarkText{La RevueCapes}
\SetWatermarkLightness{0.9}
\SetWatermarkScale{2}
\usepackage{hyperref}
\hypersetup{colorlinks=true,linkcolor=blue!50!black,urlcolor=blue!50!black}

\definecolor{revue}{RGB}{20,60,120}

\pagestyle{fancy}
\fancyhf{}
\renewcommand{\headrulewidth}{0.4pt}
\setlength{\headheight}{15pt}

% \entete{Matière}{Titre}{Type de document}
\newcommand{\entete}[3]{%
  \fancyhead[L]{\small\textsc{La Revue CAPES} -- #1}%
  \fancyhead[R]{\small #2}%
  \fancyfoot[C]{\thepage}%
  \thispagestyle{plain}%
  \begin{center}
    {\color{revue}\rule{\textwidth}{1.2pt}}\\[4pt]
    {\small\textsc{Concours directs CAP-CEG / CAPES -- Burkina Faso}}\\[6pt]
    {\LARGE\bfseries\color{revue} #2}\\[6pt]
    {\large\itshape #1 -- #3}\\[2pt]
    {\color{revue}\rule{\textwidth}{1.2pt}}
  \end{center}
  \vspace{0.5cm}%
}

% Encadré pour les remarques ajoutées dans les corrigés
\newcommand{\remarque}[1]{\par\medskip\noindent\fcolorbox{revue}{revue!6}{\parbox{\dimexpr\linewidth-2\fboxsep-2\fboxrule}{\textbf{\color{revue}Remarque.} #1}}\par\medskip}


\begin{document}
\entete{Culture générale}{Corrigé du sujet type de dissertation n°16}{Correction}
\begin{spacing}{1.5}

\noindent\textbf{Sujet.} L'orpaillage artisanal au Burkina Faso : enjeux économiques, tensions sociales et financement des groupes terroristes.

\rubrique{1. Analyse du sujet}

\begin{itemize}
\item \textbf{Type de sujet :} sujet-thème dont l'énoncé fournit lui-même les trois axes d'un \textbf{plan analytique} : I. enjeux économiques, II. tensions sociales (et environnementales), III. financement du terrorisme ; on termine par des solutions.
\item \textbf{Mots-clés :}
  \begin{itemize}
  \item « orpaillage artisanal » : exploitation de l'or par des méthodes manuelles ou semi-mécanisées, souvent informelle, sur des sites aménagés sommairement ;
  \item « enjeux économiques » : revenus, emplois, recettes de l'État, place de l'or dans l'économie (l'or est le premier produit d'exportation du Burkina Faso) ;
  \item « tensions sociales » : conflits, déscolarisation, insécurité, atteintes à la santé et à l'environnement ;
  \item « financement des groupes terroristes » : utilisation de l'or par les groupes armés pour se procurer des ressources.
  \end{itemize}
\item \textbf{Problématique :} comment l'orpaillage artisanal, source vitale de revenus pour de nombreux Burkinabè, est-il devenu aussi un facteur de tensions sociales et une ressource pour les groupes terroristes, et comment mieux l'encadrer ?
\end{itemize}

\rubrique{2. Plan détaillé}

\begin{enumerate}
\item[I.] \textbf{Des enjeux économiques considérables.} 1. Une source de revenus et d'emplois pour des centaines de milliers de personnes. 2. Un apport aux économies locales. 3. Mais un manque à gagner pour l'État (secteur informel, fraude).
\item[II.] \textbf{Des tensions sociales et environnementales.} 1. Conflits et insécurité sur les sites. 2. Déscolarisation, travail des enfants et problèmes sanitaires. 3. Dégradation de l'environnement (mercure, cyanure, destruction des terres).
\item[III.] \textbf{Une ressource pour les groupes terroristes ; quelles solutions ?} 1. Le contrôle des sites et le financement des groupes armés. 2. Les circuits de contrebande et le recrutement. 3. Mieux encadrer l'orpaillage.
\end{enumerate}

\rubrique{3. Proposition de rédaction}

\partie{Introduction}

Depuis la fin des années 2000, l'or est devenu le premier produit d'exportation du Burkina Faso, devant le coton. À côté des grandes mines industrielles, l'orpaillage artisanal occupe une place considérable : des centaines de milliers de Burkinabè, souvent jeunes, travaillent sur des sites d'orpaillage disséminés dans tout le pays. Cette activité, qui fait vivre de nombreuses familles, est aussi à l'origine de graves problèmes sociaux et environnementaux et, dans le contexte de crise sécuritaire, elle est devenue une source de financement pour les groupes terroristes. Quels sont les enjeux de l'orpaillage artisanal au Burkina Faso, et comment concilier ses avantages économiques avec la paix sociale et la sécurité ? Nous analyserons d'abord ses enjeux économiques, puis les tensions sociales qu'il génère, avant d'examiner son rôle dans le financement du terrorisme et les moyens de mieux l'encadrer.

\partie{I. Des enjeux économiques considérables}

\souspartie{1. Une source de revenus et d'emplois}
Dans un pays où l'agriculture dépend de pluies irrégulières, l'orpaillage offre une source de revenus rapide, notamment pendant la saison sèche. Il emploie directement des centaines de milliers de personnes (creuseurs, transporteurs, concasseurs, laveurs, souvent des femmes) et fait vivre indirectement de nombreuses autres. Pour beaucoup de jeunes ruraux, c'est l'une des rares alternatives au chômage et à l'exode.

\souspartie{2. Un apport aux économies locales}
Autour des sites se développent des activités commerciales : restauration, vente d'outils, de carburant, de vivres, transport, services de téléphonie. L'argent de l'or circule dans les villages, finance des constructions, des achats de bétail ou de motos, et anime l'économie locale.

\souspartie{3. Un manque à gagner pour l'État}
Mais l'essentiel de l'or artisanal échappe au contrôle de l'État : production non déclarée, ventes à des acheteurs informels, exportation frauduleuse vers les pays voisins. L'État perd ainsi des recettes fiscales importantes. Les autorités ont créé des structures pour encadrer le secteur, comme l'Agence nationale d'encadrement des exploitations minières artisanales et semi-mécanisées (ANEEMAS) et la Société de participation minière du Burkina (SOPAMIB), afin de mieux contrôler la commercialisation de l'or.

\transition{Si l'orpaillage présente des avantages économiques réels, il engendre aussi de nombreuses tensions sociales.}

\partie{II. Des tensions sociales et environnementales}

\souspartie{1. Conflits et insécurité}
La découverte d'un site attire des milliers de personnes venues d'autres régions ou de pays voisins, ce qui provoque des conflits avec les populations locales pour l'accès à la terre et aux filons, des disputes entre exploitants, des vols et des violences. Les effondrements de galeries font régulièrement des victimes.

\souspartie{2. Déscolarisation, travail des enfants et santé}
De nombreux enfants abandonnent l'école pour travailler sur les sites, dans des conditions dangereuses. La consommation de drogues, la prostitution et les maladies (infections respiratoires, VIH, accidents) y sont fréquentes. Les communautés voient leurs valeurs et leur cohésion fragilisées.

\souspartie{3. Dégradation de l'environnement}
L'orpaillage détruit les terres agricoles et les forêts, laisse des trous béants dangereux pour les hommes et le bétail, et pollue les sols et les cours d'eau par l'usage de produits chimiques comme le mercure et le cyanure, très toxiques pour la santé humaine et animale.

\transition{Dans le contexte de la crise sécuritaire, l'orpaillage a acquis une dimension encore plus préoccupante : il contribue au financement des groupes terroristes.}

\partie{III. Une ressource pour les groupes terroristes ; quelles solutions ?}

\souspartie{1. Le contrôle des sites et le financement}
Dans les zones où l'État est absent, les groupes armés ont pris le contrôle de certains sites d'orpaillage : ils prélèvent des taxes sur les orpailleurs, imposent leur « protection » ou exploitent directement l'or. Ces revenus leur permettent d'acheter des armes, des motos, du carburant et de payer leurs combattants.

\souspartie{2. Contrebande et recrutement}
L'or extrait dans ces zones rejoint des circuits de contrebande transfrontaliers difficiles à contrôler, où il est blanchi. Les sites d'orpaillage servent aussi de lieux de recrutement : des jeunes orpailleurs sont approchés, intimidés ou enrôlés par les groupes armés.

\souspartie{3. Mieux encadrer l'orpaillage}
Plusieurs mesures peuvent être envisagées : sécuriser les zones aurifères, formaliser l'activité (cartes d'orpailleur, coopératives, comptoirs agréés), renforcer la traçabilité de l'or et la lutte contre la fraude, interdire et remplacer les produits chimiques dangereux, réhabiliter les sites abandonnés, interdire le travail des enfants et offrir aux jeunes d'autres sources de revenus. La coopération régionale est indispensable pour démanteler les circuits de contrebande.

\partie{Conclusion}

L'orpaillage artisanal est une réalité ambivalente au Burkina Faso. Il représente un enjeu économique majeur, qui fait vivre de nombreuses familles et anime les économies locales, mais il échappe en grande partie au contrôle de l'État. Il génère aussi de fortes tensions sociales et de graves dégâts environnementaux, et il est devenu, dans les zones d'insécurité, une source de financement et de recrutement pour les groupes terroristes. Mieux encadrer cette activité, en la sécurisant, en la formalisant et en protégeant les populations et l'environnement, est donc un enjeu à la fois économique, social et sécuritaire. L'or du Burkina doit servir le développement du pays et non nourrir la violence.

\end{spacing}
\end{document}
